\appendix
\cleardoublepage

\chapter*{Annexes}
\addcontentsline{toc}{chapter}{Annexes}
\markboth{Annexes}{Annexes}

\section*{Annexe A : Commandes techniques de la plateforme}
\addcontentsline{toc}{section}{Annexe A : Commandes techniques de la plateforme}
\label{annexe:commandes-techniques}

Cette annexe regroupe les commandes essentielles utilisées pour le développement, la gestion du code source, la conteneurisation et le déploiement de la plateforme \textit{AiContentFlow}. Elle sert de référence technique pour reproduire les principales opérations réalisées durant le projet et faciliter l'onboarding de nouveaux développeurs.

\vspace{0.5cm}

\textbf{Vérification de l'environnement .NET}

\begin{verbatim}
dotnet --version
\end{verbatim}

\vspace{0.3cm}

\textbf{Initialisation du dépôt Git}

\begin{verbatim}
git init
git clone https://github.com/AYGTX/aicontentflow.git
git status
\end{verbatim}

\vspace{0.3cm}

\textbf{Configuration des variables d'environnement}

\begin{verbatim}
VITE_API_URL=http://localhost:5000
VITE_ENVIRONMENT=development
\end{verbatim}

\vspace{0.3cm}

\textbf{Fichier .gitignore}

\begin{verbatim}
node_modules/
dist/
build/
.env
.env.local
bin/
obj/
.vscode/
.idea/
*.log
\end{verbatim}

\vspace{0.5cm}

\textbf{Création et structure de la solution backend ASP.NET Core}

\begin{verbatim}
dotnet new sln -n AiContentFlow
dotnet new classlib -n AiContentFlow.Domain
dotnet new classlib -n AiContentFlow.Application
dotnet new classlib -n AiContentFlow.Infrastructure
dotnet new webapi -n AiContentFlow.API
dotnet sln add AiContentFlow.Domain/AiContentFlow.Domain.csproj
dotnet sln add AiContentFlow.Application/AiContentFlow.Application.csproj
dotnet sln add AiContentFlow.Infrastructure/AiContentFlow.Infrastructure.csproj
dotnet sln add AiContentFlow.API/AiContentFlow.API.csproj
\end{verbatim}

\vspace{0.3cm}

\textbf{Restauration des dépendances backend}

\begin{verbatim}
dotnet restore
\end{verbatim}

\vspace{0.3cm}

\textbf{Développement et exécution du backend}

\begin{verbatim}
dotnet run
dotnet watch run
dotnet build
dotnet build -c Release
\end{verbatim}

\vspace{0.3cm}

\textbf{Gestion de la base de données avec Entity Framework Core}

\begin{verbatim}
dotnet tool install --global dotnet-ef
dotnet ef migrations add InitialCreate
dotnet ef database update
dotnet ef migrations list
dotnet ef migrations remove
\end{verbatim}

\vspace{0.5cm}

\textbf{Création du projet frontend React avec Vite}

\begin{verbatim}
npm create vite@latest aicontentflow-web -- --template react-ts
cd aicontentflow-web
npm install
\end{verbatim}

\vspace{0.3cm}

\textbf{Exécution et compilation du frontend}

\begin{verbatim}
npm run dev
npm run build
npm run preview
npm run lint
\end{verbatim}

\vspace{0.3cm}

\textbf{Installation des dépendances du frontend}

\begin{verbatim}
npm install react-router-dom
npm install axios
npm install @tanstack/react-query
npm install zustand
npm install react-hook-form
npm install zod
npm install @mui/material
npm install framer-motion
\end{verbatim}

\vspace{0.5cm}

\textbf{Opérations Git de base}

\begin{verbatim}
git status
git add .
git commit -m "Description des modifications"
git push origin main
git pull origin main
\end{verbatim}

\vspace{0.3cm}

\textbf{Gestion des branches Git}

\begin{verbatim}
git branch
git checkout -b feature/nouvelle-fonctionnalite
git merge feature/nouvelle-fonctionnalite
git branch -d feature/nouvelle-fonctionnalite
\end{verbatim}

\vspace{0.5cm}

\textbf{Administration PostgreSQL}

\begin{verbatim}
psql -U postgres
CREATE DATABASE aicontentflow_db;
\c aicontentflow_db
\dt
pg_dump -U postgres aicontentflow_db > backup.sql
psql -U postgres aicontentflow_db < backup.sql
\q
\end{verbatim}

\vspace{0.5cm}

\textbf{Construction des images Docker}

\begin{verbatim}
docker build -t aicontentflow-api .
docker build -f Dockerfile.frontend -t aicontentflow-web .
\end{verbatim}

\vspace{0.3cm}

\textbf{Gestion des conteneurs Docker}

\begin{verbatim}
docker run -d -p 5000:5000 aicontentflow-api
docker ps
docker ps -a
docker stop <container_id>
docker restart <container_id>
docker logs <container_id>
docker logs -f <container_id>
\end{verbatim}

\vspace{0.3cm}

\textbf{Docker Compose}

\begin{verbatim}
docker-compose up -d
docker-compose down
docker-compose up -d --build
docker-compose logs -f
\end{verbatim}

\vspace{0.5cm}

\textbf{Cycle de développement du backend}

\begin{verbatim}
dotnet restore
dotnet watch run
git add .
git commit -m "feat: description de la feature"
git push origin feature-branch
dotnet build -c Release
\end{verbatim}

\vspace{0.3cm}

\textbf{Cycle de développement du frontend}

\begin{verbatim}
npm install
npm run dev
npm run lint
npm run build
git add .
git commit -m "feat: description de la feature"
git push origin feature-branch
\end{verbatim}

\vspace{0.3cm}

\textbf{Déploiement complet}

\begin{verbatim}
docker-compose down
docker-compose up -d --build
docker-compose ps
docker-compose logs -f
pg_dump -U postgres aicontentflow_db > backup_$(date +%Y%m%d_%H%M%S).sql
\end{verbatim}

\vspace{0.3cm}

\textbf{Résolution de problèmes courants}

\begin{verbatim}
dotnet clean
dotnet restore
dotnet build
rm -r node_modules package-lock.json
npm install
docker-compose down -v
docker-compose up -d --build
dotnet run 2>&1 | grep -i error
psql -U postgres -h localhost -d aicontentflow_db -c "SELECT VERSION();"
\end{verbatim}
